% !TeX program = pdflatex
% Compile twice: pdflatex pv_eta_invariant_en.tex
% Self-contained source: bibliography included; no external figures required.
% Author and affiliation can be supplied in the author field below.
\documentclass[11pt,a4paper]{article}
\usepackage[T1]{fontenc}
\usepackage{lmodern}
\usepackage[margin=27mm,headheight=14pt]{geometry}
\usepackage{amsmath,amssymb,amsthm,mathtools}
\usepackage{microtype,booktabs,fancyhdr,xcolor}
\usepackage[colorlinks=true,linkcolor=blue!40!black,citecolor=blue!40!black,urlcolor=blue!40!black]{hyperref}
\usepackage{bookmark}
\hypersetup{pdftitle={Finite Pauli--Villars Products and the Eta Invariant},
  pdfsubject={Convergent spectral determinants, Pfaffians, and large-mass limits}}
\newtheorem{theorem}{Theorem}[section]
\newtheorem{proposition}[theorem]{Proposition}
\newtheorem{lemma}[theorem]{Lemma}
\newtheorem{corollary}[theorem]{Corollary}
\theoremstyle{remark}
\newtheorem{remark}[theorem]{Remark}
\DeclareMathOperator{\Tr}{Tr}
\DeclareMathOperator{\Pf}{Pf}
\DeclareMathOperator{\sgn}{sgn}
\DeclareMathOperator{\spec}{spec}
\DeclareMathOperator{\Log}{Log}
\DeclareMathOperator{\ph}{ph}
\newcommand{\ii}{\mathrm i}
\newcommand{\dd}{\,\mathrm d}
\newcommand{\C}{\mathbb C}
\newcommand{\R}{\mathbb R}
\newcommand{\Maj}{\mathrm{Maj}}
\numberwithin{equation}{section}
\allowdisplaybreaks[2]
\emergencystretch=2em
\setlength{\parskip}{2pt}
\pagestyle{fancy}
\fancyhf{}
\fancyhead[L]{\small Finite Pauli--Villars Products and the Eta Invariant}
\fancyfoot[C]{\thepage}
\renewcommand{\headrulewidth}{0.3pt}

\title{\bfseries Finite Pauli--Villars Products\\and the Eta Invariant}
\author{}
\date{October 8, 2026}

\begin{document}
\maketitle

\begin{abstract}
We give an explicit spectral construction of Pauli--Villars regularization
for geometric Dirac operators on closed odd-dimensional spin manifolds.
A finite collection of regulator masses, with binomial multiplicities,
defines an ordinary absolutely convergent product at every positive
regulator scale. Using the Bismut--Freed local heat-kernel estimate, we
prove that its phase converges to $\exp(-\ii\pi\eta_D(0)/2)$.
For this particular scheme, the modulus also has a finite limit, equal to
the positive square root of the zeta determinant of $D^2$ when $D$ is
invertible. An antiunitary symmetry with square $-1$ gives a compatible
Pfaffian construction and the phase $\exp(-\ii\pi\eta_D(0)/4)$.
We also treat negative physical mass and zero modes, and isolate the
analytic hypotheses needed beyond the odd-dimensional oriented setting.
The purpose is to provide a direct convergence proof of the established
eta-phase formula, with the regulator scheme and the Pfaffian branch
specified explicitly.
\end{abstract}

\section{Introduction}\label{sec:intro}

The phase of a fermion partition function contains information that is
not determined by its absolute value. For a massless complex fermion,
a familiar Pauli--Villars (PV) prescription leads to the expression
\begin{equation}\label{eq:intro-phase}
 Z=|Z|\exp\!\left(-\frac{\ii\pi}{2}\eta_D(0)\right),
\end{equation}
where $D$ is a self-adjoint Dirac operator and the sign depends on the
regulator-mass convention. This formula, and its Pfaffian counterpart,
play a central role in the discussion of fermion anomalies and
topological phases in \cite{Witten,WY}.

To make a statement such as \eqref{eq:intro-phase} precise, one must
specify both the regularized determinant and the limiting operation.
For example, the formal product
$\prod_\lambda\lambda/(\lambda+\ii M)$ is not an ordinary absolutely
convergent product in three dimensions. Taking a large-mass limit
eigenvalue by eigenvalue does not justify interchanging that limit with
an infinite spectral sum. Moreover, a Pfaffian requires a compatible
choice of sign, rather than an unspecified square root of a determinant.

We address these points by an explicit finite PV construction. For any
even integer $N$ larger than the dimension, the coefficients
\begin{equation}\label{eq:intro-c}
 c_j=(-1)^j\binom Nj,\qquad 0\le j\le N,
\end{equation}
cancel all moments of degree less than $N$. Grouping the regulator
factors at each eigenvalue then produces an ordinary Fredholm
determinant. Its phase admits a convergent heat-kernel representation.
The Bismut--Freed estimate for the local eta density supplies an
integrable majorant, so that the large-mass limit follows from dominated
convergence. The same moment conditions also cancel the divergent
terms in the modulus in odd dimensions.

The analytic input is the local regularity theorem of Bismut and
Freed \cite{BF}; an alternative proof is given by Ponge \cite{Ponge}.
The geometry of determinant and Pfaffian lines, including families and
boundary problems, is developed in \cite{DF,FM}. The present paper
concerns explicit scalar spectral products for a fixed operator and
does not replace that geometric framework. The phase formula itself
is established in the literature. Our aim is to exhibit a concrete
finite-regulator realization and to justify each interchange of a sum,
an integral, and a large-mass limit.

Section~\ref{sec:setup} records the geometric and analytic assumptions.
Section~\ref{sec:main} states and proves the determinant theorem.
Section~\ref{sec:pf} gives the Pfaffian construction.
Negative mass and zero modes are treated in Section~\ref{sec:mass}.
Section~\ref{sec:extensions} discusses extensions and limitations.

\section{Geometric setting and analytic input}\label{sec:setup}

Let $Y$ be a closed Riemannian spin manifold of odd dimension
$n=2r+1$, and let $E\to Y$ be a Hermitian vector bundle with a unitary
connection. Write $S\to Y$ for the complex spinor bundle. We consider
the geometric twisted Dirac operator
\[
 D:H^1(Y;S\otimes E)\longrightarrow L^2(Y;S\otimes E)
\]
as a self-adjoint unbounded operator on $L^2(Y;S\otimes E)$.
Here ``geometric'' refers to the operator formed from the Levi--Civita
and twisting connections, without an additional arbitrary zero-order
potential. The regulator convention throughout is
\begin{equation}\label{eq:mass-sign}
 D+\ii m,\qquad m>0.
\end{equation}
Replacing every regulator mass by its negative conjugates the
massless phase obtained below.

Spectral sums count eigenvalues with their complex multiplicities.
In particular, the independent Grassmann variables of a complex
fermion and its conjugate do not introduce a second copy of this
spectral multiplicity. When a kernel is present, zero eigenvalues are
omitted from the eta function:
\begin{equation}\label{eq:eta-def}
 \eta_D(s)=\sum_{\lambda\ne0}\sgn(\lambda)|\lambda|^{-s},
 \qquad \Re s>n.
\end{equation}
We set $h=\dim_\C\ker D$ and use $\ph z=z/|z|$ for $z\ne0$.

Introduce the heat traces
\begin{equation}\label{eq:heat-traces}
 K(t)=\Tr(e^{-tD^2}),\qquad
 T(t)=\Tr(De^{-tD^2}),\qquad
 A(t)=\Tr(|D|e^{-tD^2}).
\end{equation}
The Weyl estimate implies
\begin{equation}\label{eq:weyl}
 \#\{\lambda:0<|\lambda|\le R\}=O(R^n),\qquad
 A(t)=O(t^{-(n+1)/2})\quad(t\downarrow0).
\end{equation}
Both $T(t)$ and $A(t)$ decay exponentially as $t\to\infty$ because
the nonzero spectrum has a positive gap from zero. The same is true
of $K(t)$ when $h=0$.

\begin{theorem}[Bismut--Freed local regularity]\label{thm:BF}
For the geometric Dirac operator above, the fiberwise trace of the
kernel of $De^{-tD^2}$ on the diagonal is uniformly $O(t^{1/2})$ as
$t\downarrow0$. Consequently,
\begin{equation}\label{eq:BF}
 T(t)=O(t^{1/2}),\qquad
 \int_0^\infty t^{-1/2}|T(t)|\dd t<\infty.
\end{equation}
\end{theorem}

This is the analytic result of \cite[Theorem 2.4]{BF}; see also
\cite[Theorem 3.1]{Ponge}. We use it as an established theorem.

\begin{lemma}\label{lem:eta}
The eta function is holomorphic for $\Re s>-2$ and satisfies
\begin{equation}\label{eq:eta-heat}
 \eta_D(0)=\frac1{\sqrt\pi}\int_0^\infty t^{-1/2}T(t)\dd t,
\end{equation}
where the integral is absolutely convergent.
\end{lemma}
\begin{proof}
For $\Re s>n$, absolute convergence permits the Mellin representation
\begin{equation}\label{eq:mellin}
 \eta_D(s)=\frac1{\Gamma((s+1)/2)}
 \int_0^\infty t^{(s-1)/2}T(t)\dd t.
\end{equation}
Theorem~\ref{thm:BF} and exponential decay at infinity show that the
integral is locally uniformly convergent for $\Re s>-2$.
It therefore gives a holomorphic continuation to that half-plane.
Evaluation at $s=0$ gives \eqref{eq:eta-heat}.
\end{proof}

The argument below needs the absolute integrability in
\eqref{eq:BF}, not merely regularity of an abstract eta function at
$s=0$. This distinction matters when considering more general
self-adjoint elliptic operators.

\section{The finite PV determinant}\label{sec:main}

Assume first that $h=0$. Choose an even integer $N>n$ and take the
coefficients \eqref{eq:intro-c}. For $M>0$ and $\lambda\ne0$, set
\begin{equation}\label{eq:g-def}
 g_M(\lambda)=\prod_{j=0}^N(\lambda+\ii jM)^{c_j},\qquad
 Z_M(D)=\prod_{\lambda\in\spec D}g_M(\lambda).
\end{equation}
The factor with $j=0$ is the physical massless determinant. Positive
and negative coefficients specify positive and inverse determinant
powers, respectively. The factors are combined at each eigenvalue
before the infinite product is taken.

Write the short-time expansion of the heat trace as
\begin{equation}\label{eq:K-expansion}
 K(t)=\sum_{k=0}^r a_{2k}t^{k-n/2}+O(t^{1/2}),
 \qquad
 Q(t)=K(t)-\sum_{k=0}^r a_{2k}t^{k-n/2}.
\end{equation}
Define
\begin{equation}\label{eq:amplitude}
 \mathcal A(D)=
 \exp\!\left[-\frac12\int_0^\infty Q(t)\frac{\dd t}{t}\right].
\end{equation}
The integral converges at both endpoints: $Q(t)/t=O(t^{-1/2})$
near zero, while $K(t)$ decays exponentially and the subtracted
powers decay at infinity.

\begin{theorem}[Determinant limit]\label{thm:main}
For every $M>0$, the product \eqref{eq:g-def} is absolutely convergent
and nonzero. Moreover, $g_M(D)-1$ is trace class and
\begin{equation}\label{eq:fredholm}
 Z_M(D)=\det\nolimits_F g_M(D).
\end{equation}
Its large-mass limit is
\begin{equation}\label{eq:main-limit}
 \lim_{M\to\infty}Z_M(D)
 =\mathcal A(D)\exp\!\left(-\frac{\ii\pi}{2}\eta_D(0)\right).
\end{equation}
In particular,
\begin{equation}\label{eq:phase-limit}
 \lim_{M\to\infty}\ph Z_M(D)
 =\exp\!\left(-\frac{\ii\pi}{2}\eta_D(0)\right).
\end{equation}
Both limits are independent of the even integer $N>n$ within this
family of regulator schemes.
\end{theorem}

\begin{corollary}[Three-dimensional formula]\label{cor:3d}
If $n=3$, one may take $N=4$ and obtain
\begin{equation}\label{eq:3d-product}
 Z_M(D)=\prod_{\lambda\in\spec D}
 \frac{\lambda(\lambda+2\ii M)^6(\lambda+4\ii M)}
      {(\lambda+\ii M)^4(\lambda+3\ii M)^4}.
\end{equation}
This product satisfies \eqref{eq:main-limit} with
\begin{equation}\label{eq:3d-amplitude}
 \mathcal A(D)=\exp\!\left[-\frac12\int_0^\infty
 \frac{K(t)-a_0t^{-3/2}-a_2t^{-1/2}}{t}\dd t\right].
\end{equation}
\end{corollary}

The regulator data in this example are
\[
\begin{array}{c|rrrrr}
 j&0&1&2&3&4\\\hline
 c_j&1&-4&6&-4&1\\
 m_j&0&M&2M&3M&4M
\end{array}
\]
Only the combined product is used. We do not define the individual
divergent determinants and then assume multiplicativity of their
regularized values.

\subsection{Absolute convergence at finite mass}

The finite-difference identity gives
\begin{equation}\label{eq:moments}
 \sum_{j=0}^N c_jp(j)=0
 \quad\text{for every polynomial $p$ of degree less than $N$}.
\end{equation}
In particular, $\sum c_j=0$ and $\sum c_jj^k=0$ for $1\le k<N$.
Fix the logarithm
\begin{equation}\label{eq:ell}
 \ell_M(\lambda)=\sum_{j=1}^Nc_j
       \Log\!\left(1+\frac{\ii jM}{\lambda}\right),
\end{equation}
where every $\Log$ is the principal logarithm. Its argument has real
part one, so no factor meets the branch cut. Since $\sum c_j=0$,
$e^{\ell_M(\lambda)}=g_M(\lambda)$ for either sign of $\lambda$.

For fixed $M$, expansion at large $|\lambda|$ and
\eqref{eq:moments} give
\begin{equation}\label{eq:tail}
 \ell_M(\lambda)=O_M(|\lambda|^{-N}).
\end{equation}
The Weyl estimate and $N>n$ imply
\begin{equation}\label{eq:abs-sum}
 \sum_\lambda|\ell_M(\lambda)|<\infty.
\end{equation}
It follows that the unordered product is well defined and nonzero,
with
\begin{equation}\label{eq:Z-log}
 Z_M(D)=\exp\!\left(\sum_\lambda\ell_M(\lambda)\right).
\end{equation}
Also $\sum_\lambda|g_M(\lambda)-1|<\infty$. In a spectral basis this
is precisely the trace-class condition for $g_M(D)-1$, and it proves
\eqref{eq:fredholm}.

\subsection{A convergent formula for the phase}

The logarithm \eqref{eq:ell} fixes a real-valued lift of the phase:
\begin{equation}\label{eq:theta}
 \Theta_M=\operatorname{Im}\sum_\lambda\ell_M(\lambda)
 =\sum_\lambda\sum_{j=1}^Nc_j\arctan\frac{jM}{\lambda}.
\end{equation}
The arctangent takes values in $(-\pi/2,\pi/2)$.
For $a\in\R$ and $\lambda\ne0$,
\begin{align}\label{eq:arctan}
 \arctan\frac a\lambda
 &=\int_0^a\frac{\lambda}{\lambda^2+u^2}\dd u\notag\\
 &=\int_0^\infty\lambda e^{-t\lambda^2}
            \left(\int_0^a e^{-tu^2}\dd u\right)\dd t.
\end{align}
The inner integral is oriented when $a<0$.
Introduce
\begin{equation}\label{eq:H}
 H_M(t)=\sum_{j=1}^Nc_j\int_0^{jM}e^{-tu^2}\dd u.
\end{equation}
At fixed $M$, its expansion at zero is
\begin{equation}\label{eq:H-expansion}
 H_M(t)=\sum_{q\ge0}
 \frac{(-t)^qM^{2q+1}}{q!(2q+1)}\sum_{j=1}^Nc_jj^{2q+1}.
\end{equation}
Since $N$ is even, the terms with $0\le q<N/2$ vanish. Hence
\begin{equation}\label{eq:H-small}
 H_M(t)=O_M(t^{N/2}).
\end{equation}
Together with \eqref{eq:weyl}, this gives
\begin{equation}\label{eq:fubini}
 \int_0^\infty\sum_\lambda
 |H_M(t)\lambda e^{-t\lambda^2}|\dd t
 =\int_0^\infty |H_M(t)|A(t)\dd t<\infty.
\end{equation}
Indeed, the integrand is $O_M(t^{(N-n-1)/2})$ at zero and is
integrable at infinity. Thus Fubini's theorem applies after the PV
terms have been grouped in $H_M$, and yields
\begin{equation}\label{eq:theta-heat}
 \Theta_M=\int_0^\infty H_M(t)T(t)\dd t.
\end{equation}

For fixed $t>0$, $\sum_{j=1}^Nc_j=-1$ implies
\begin{equation}\label{eq:H-limit}
 H_M(t)\longrightarrow-\frac{\sqrt\pi}{2\sqrt t}
 \quad(M\to\infty).
\end{equation}
Uniformly in $M$,
\begin{equation}\label{eq:H-majorant}
 |H_M(t)|\le(2^N-1)\frac{\sqrt\pi}{2\sqrt t}.
\end{equation}
The integrability in \eqref{eq:BF} therefore permits dominated
convergence in \eqref{eq:theta-heat}. By Lemma~\ref{lem:eta},
\begin{equation}\label{eq:theta-limit}
 \lim_{M\to\infty}\Theta_M
 =-\frac{\sqrt\pi}{2}\int_0^\infty t^{-1/2}T(t)\dd t
 =-\frac\pi2\eta_D(0).
\end{equation}
This proves the phase statement. In fact, the chosen real lift itself
converges; the conclusion is stronger than convergence modulo $2\pi$.
Heat kernels have only been used to analyze the already defined PV
product, rather than as an additional prescription for defining it.

\subsection{The modulus and cancellation of local terms}

Set
\begin{equation}\label{eq:B}
 B_M(t)=\sum_{j=0}^Nc_j e^{-j^2M^2t}.
\end{equation}
The even moments in \eqref{eq:moments} give
$B_M(t)=O_M(t^{N/2})$ at zero. For each nonzero eigenvalue, the
Frullani identity and $\sum c_j=0$ imply
\begin{equation}\label{eq:frullani}
 \sum_{j=0}^Nc_j\log(\lambda^2+j^2M^2)
 =-\int_0^\infty e^{-t\lambda^2}B_M(t)\frac{\dd t}{t}.
\end{equation}
The summed absolute integrand is integrable because
\[
 |B_M(t)|K(t)/t=O_M(t^{(N-n)/2-1})\quad(t\downarrow0),
\]
and $K(t)$ decays exponentially at infinity. Consequently,
\begin{equation}\label{eq:log-modulus}
 \log|Z_M(D)|=-\frac12\int_0^\infty B_M(t)K(t)\frac{\dd t}{t}.
\end{equation}

\begin{lemma}[Vanishing heat moments]\label{lem:heat-moments}
For every odd integer $p$ with $1\le p\le n$,
\begin{equation}\label{eq:heat-moments}
 \int_0^\infty t^{-1-p/2}B_M(t)\dd t
 =\Gamma(-p/2)M^p\sum_{j=1}^Nc_jj^p=0.
\end{equation}
All integrals in the left-hand expression are ordinary convergent
integrals of the combined function $B_M$.
\end{lemma}
\begin{proof}
Convergence follows from $N>p$ at zero and boundedness of $B_M$ at
infinity. Let $q=(p+1)/2$. Integrating by parts $q$ times gives
\begin{equation}\label{eq:ibp}
 \frac{(-1)^q}{\prod_{\ell=0}^{q-1}(-p/2+\ell)}
 \int_0^\infty t^{q-1-p/2}B_M^{(q)}(t)\dd t.
\end{equation}
At each step the boundary terms vanish: at zero one uses the
vanishing order of $B_M$ and its derivatives; at infinity the initial
power is negative and all subsequent derivatives decay exponentially.
After differentiating, the $j=0$ term is absent and each remaining
exponential is individually integrable. Evaluating these integrals
and using the recurrence of the gamma function gives the middle
expression in \eqref{eq:heat-moments}. The odd moment in
\eqref{eq:moments} then makes it zero.
\end{proof}

Each subtracted heat term in \eqref{eq:K-expansion} corresponds to
$p=n-2k$, a positive odd integer. Lemma~\ref{lem:heat-moments} thus
turns \eqref{eq:log-modulus} into the exact identity
\begin{equation}\label{eq:modulus-Q}
 \log|Z_M(D)|=-\frac12\int_0^\infty B_M(t)Q(t)\frac{\dd t}{t}.
\end{equation}
Since $|B_M(t)|\le2^N$, $B_M(t)\to1$ for every $t>0$, and
$Q(t)/t$ is absolutely integrable, dominated convergence yields
\begin{equation}\label{eq:modulus-limit}
 \lim_{M\to\infty}|Z_M(D)|=\mathcal A(D)>0.
\end{equation}
Combining this with \eqref{eq:theta-limit} proves
Theorem~\ref{thm:main} and Corollary~\ref{cor:3d}.

In three dimensions, the two cancellations read
\begin{align}
 \int_0^\infty t^{-3/2}B_M(t)\dd t
 &=-2\sqrt\pi M\sum_{j=1}^4c_jj=0,\label{eq:3d-moment1}\\
 \int_0^\infty t^{-5/2}B_M(t)\dd t
 &=\frac{4\sqrt\pi}{3}M^3\sum_{j=1}^4c_jj^3=0.
 \label{eq:3d-moment3}
\end{align}
Thus the subtractions in the proof are exact zero identities at
finite $M$. No divergent real counterterm has been added to this
particular determinant prescription.

\subsection{Identification of the limiting modulus}

For comparison with a conventional normalization, define
\[
 \zeta_{D^2}(s)=\sum_\lambda(\lambda^2)^{-s},
 \qquad \det\nolimits_\zeta D^2=\exp[-\zeta_{D^2}'(0)].
\]

\begin{proposition}\label{prop:zeta}
For invertible $D$ in the setting of Theorem~\ref{thm:main},
\begin{equation}\label{eq:zeta-modulus}
 \mathcal A(D)=(\det\nolimits_\zeta D^2)^{1/2}.
\end{equation}
\end{proposition}
\begin{proof}
Put $p_k=n-2k$ for $0\le k\le r$. The meromorphic continuation of
the Mellin transform is, near $s=0$,
\begin{align}
 \Gamma(s)\zeta_{D^2}(s)
 &=\int_0^1 t^{s-1}Q(t)\dd t
   +\int_1^\infty t^{s-1}K(t)\dd t
   +\sum_{k=0}^r\frac{a_{2k}}{s-p_k/2}\notag\\
 &=\int_0^\infty t^{s-1}Q(t)\dd t.\label{eq:zeta-Q}
\end{align}
The second equality follows by integrating the subtracted powers
over $[1,\infty)$. The last integral is holomorphic in the strip
$-1/2<\Re s<1/2$. Since $1/\Gamma(s)=s+O(s^2)$,
\[
 \zeta_{D^2}'(0)=\int_0^\infty Q(t)\frac{\dd t}{t}.
\]
Equation~\eqref{eq:zeta-modulus} follows from \eqref{eq:amplitude}.
\end{proof}

This comparison is a consequence of the convergent construction;
zeta regularization was not used to define the finite PV product.

\section{Pfaffians and Majorana fermions}\label{sec:pf}

A Pfaffian formula requires additional geometric structure. We give
a sufficient condition that specifies both the antisymmetric form
and a compatible square root of the determinant.

\begin{theorem}[Pfaffian limit]\label{thm:pf}
Assume the hypotheses of Theorem~\ref{thm:main}. Suppose that a
fiberwise antilinear isometry induces an antiunitary operator $J$ on
the Hilbert space, preserves the domain of $D$, and satisfies
\begin{equation}\label{eq:J}
 J^2=-1,\qquad JD=DJ.
\end{equation}
Take the Hermitian inner product to be linear in its second argument
and define
\begin{equation}\label{eq:omega}
 \omega(u,v)=\langle Ju,v\rangle,
 \qquad A_m(u,v)=\omega\bigl(u,(D+\ii m)v\bigr).
\end{equation}
Then $A_m$ is a skew-symmetric complex bilinear form. The finite
spectral Pfaffian ratios with coefficients $c_j$ have an absolutely
convergent limit $P_M(D)$ as the spectral cutoff is removed, with
\begin{equation}\label{eq:pf-square}
 P_M(D)^2=Z_M(D).
\end{equation}
Furthermore,
\begin{equation}\label{eq:pf-limit}
 \lim_{M\to\infty}P_M(D)
 =\mathcal A(D)^{1/2}
       \exp\!\left(-\frac{\ii\pi}{4}\eta_D(0)\right).
\end{equation}
The eta invariant here counts the full complex multiplicity,
including each pair enforced by $J$.
\end{theorem}

\begin{proof}
Antiunitarity and $J^2=-1$ imply that $\omega$ is nondegenerate,
complex bilinear, and skew-symmetric. Self-adjointness and
$JD=DJ$ give
\begin{equation}\label{eq:omega-D}
 \omega(Du,v)=\langle JDu,v\rangle
 =\langle DJu,v\rangle=\langle Ju,Dv\rangle=\omega(u,Dv).
\end{equation}
Hence $A_m$ is also skew-symmetric.

Let $V_\lambda$ be the eigenspace associated with a distinct real
eigenvalue $\lambda$. It is preserved by $J$ and has even complex
dimension, say $\dim_\C V_\lambda=2r_\lambda$.
The restriction of $\omega$ to $V_\lambda$ is nondegenerate, and
\begin{equation}\label{eq:pf-eigenspace}
 \Pf(A_m|_{V_\lambda})
 =(\lambda+\ii m)^{r_\lambda}\Pf(\omega|_{V_\lambda}).
\end{equation}
Use the same basis for every mass on this eigenspace. In the product
with exponents $c_j$, the factor $\Pf(\omega|_{V_\lambda})$ cancels
because $\sum c_j=0$. Basis-change factors cancel for the same reason.
Thus the regularized Pfaffian is
\begin{equation}\label{eq:pf-product}
 P_M(D)=\prod_{\lambda\ \mathrm{distinct}}g_M(\lambda)^{r_\lambda}
 =\exp\!\left(\frac12\sum_{\lambda\ \mathrm{with\ multiplicity}}
                         \ell_M(\lambda)\right).
\end{equation}
Absolute convergence follows from \eqref{eq:abs-sum}.
This proves \eqref{eq:pf-square} and fixes its square-root branch
through finite-dimensional Pfaffian ratios. Its lifted phase is
$\Theta_M/2$. Equations~\eqref{eq:theta-limit} and
\eqref{eq:modulus-limit} prove \eqref{eq:pf-limit}.
\end{proof}

For example, the usual spin Dirac operator in dimension $8k+3$,
twisted by a bundle with a compatible real structure, has the
quaternionic symmetry relevant to this construction; see
\cite[Section 1.3]{FM}. The existence of a Majorana field with the
required scalar mass term depends on the dimension and gauge
representation. Condition~\eqref{eq:J} is sufficient for the theorem
and is not a classification of all real fermion theories.

There are two useful checks on the construction. First, when a single
$J$-pair of eigenvalues crosses zero from negative to positive, the
discrete jump in $\eta_D(0)$ is four. The phase in
\eqref{eq:pf-limit} changes sign, as a Pfaffian should. A pointwise
principal square root of the complex determinant would not in general
retain this information.

Second, one cannot obtain an inverse Pfaffian by putting a
skew-symmetric quadratic form on a single set of commuting variables.
In our scheme every negative coefficient is even. Indeed, $c_j<0$
requires $j$ odd, while
\[
 j\binom Nj=N\binom{N-1}{j-1}
\]
shows that $\binom Nj$ is even when $N$ is even and $j$ is odd.
The inverse Pfaffians can therefore be paired into inverse
determinants. The spectral definition itself is already complete
without a separate formal functional-integral representation.

\section{Negative physical mass and zero modes}\label{sec:mass}

\subsection{The negative-mass phase}

The large negative-mass phase has twice the exponent of the
massless formula. This can be proved by a shifted version of the
same finite-difference scheme. We now allow $h=\dim_\C\ker D$ to
be nonzero, retaining the convention \eqref{eq:eta-def} for eta.

\begin{theorem}[Negative physical mass]\label{thm:negative}
Let $D$ be as in Section~\ref{sec:setup}, let $N>n$ be even, and
take $c_j$ as in \eqref{eq:intro-c}. Define
\begin{equation}\label{eq:negative-product}
 a_j=2j-1,
 \qquad W_M(D)=\prod_{\lambda\in\spec D}
                       \prod_{j=0}^N(\lambda+\ii a_jM)^{c_j}.
\end{equation}
Thus the physical mass is $-M$ and all regulator masses are positive.
For every $M>0$, this is an absolutely convergent nonzero product,
including its factors at zero eigenvalues, and
\begin{equation}\label{eq:negative-phase}
 \lim_{M\to\infty}\ph W_M(D)
 =\exp[-\ii\pi(\eta_D(0)+h)].
\end{equation}
If a symmetry $J$ satisfying \eqref{eq:J} exists, the corresponding
Pfaffian ratio obeys
\begin{equation}\label{eq:negative-pf}
 \lim_{M\to\infty}\ph W_M^{\Maj}(D)
 =\exp\!\left[-\frac{\ii\pi}{2}(\eta_D(0)+h)\right].
\end{equation}
\end{theorem}

\begin{proof}
Since $a_j$ is affine in $j$, $\sum c_ja_j^k=0$ for $0\le k<N$.
For $\lambda\ne0$, the logarithm
\[
 \sum_{j=0}^Nc_j\Log(1+\ii a_jM/\lambda)
\]
is $O_M(|\lambda|^{-N})$. This proves convergence of the nonzero
spectral product. Its phase is handled by replacing $H_M$ with
\begin{equation}\label{eq:H-negative}
 \widetilde H_M(t)=\sum_{j=0}^Nc_j
                         \int_0^{a_jM}e^{-tu^2}\dd u.
\end{equation}
The same moment cancellations give
$\widetilde H_M(t)=O_M(t^{N/2})$ at zero, while
\begin{equation}
 |\widetilde H_M(t)|\le 2^N\frac{\sqrt\pi}{2\sqrt t}.
\end{equation}
As $a_0=-1$ and $a_j>0$ for $j\ge1$,
\begin{equation}\label{eq:H-negative-limit}
 \lim_{M\to\infty}\widetilde H_M(t)
 =\frac{\sqrt\pi}{2\sqrt t}\sum_{j=0}^Nc_j\sgn(a_j)
 =-\frac{\sqrt\pi}{\sqrt t}.
\end{equation}
Fubini's theorem and dominated convergence apply exactly as before.
The nonzero modes contribute $\exp(-\ii\pi\eta_D(0))$ for the
determinant and $\exp(-\ii\pi\eta_D(0)/2)$ for the Pfaffian.

A zero eigenvalue contributes the finite, strictly negative factor
\begin{equation}\label{eq:zero-factor}
 \prod_{j=0}^N(\ii a_jM)^{c_j}
 =-\prod_{j=1}^Na_j^{c_j}.
\end{equation}
The determinant therefore receives an additional sign $(-1)^h$.
In the Pfaffian case $h$ is even, and the corresponding sign is
$(-1)^{h/2}$. These signs give \eqref{eq:negative-phase} and
\eqref{eq:negative-pf}.
\end{proof}

In terms of the reduced eta invariant
$\bar\eta_D=(\eta_D(0)+h)/2$, these phases are
$e^{-2\pi\ii\bar\eta_D}$ and $e^{-\pi\ii\bar\eta_D}$, respectively.
Here $\bar\eta_D$ denotes the specified real number; in the Pfaffian
expression, reducing it modulo integers would discard a sign.
These statements concern the phase. The modulus of this massive
product generally requires a normalization by real local terms.
The cancellation proved for the massless modulus in
Section~\ref{sec:main} does not assert a corresponding massive limit.

\subsection{Massless determinants with the kernel omitted}

If $h>0$, the massless determinant itself vanishes and has no phase.
One may nevertheless define $Z'_M(D)$ by the product
\eqref{eq:g-def} over nonzero eigenvalues only. All phase arguments
remain valid, since zero modes do not contribute to $T(t)$.

\begin{proposition}\label{prop:primed}
For the massless primed product,
\begin{equation}\label{eq:primed-phase}
 \lim_{M\to\infty}\ph Z'_M(D)
 =\exp\!\left(-\frac{\ii\pi}{2}\eta_D(0)\right).
\end{equation}
There is a finite positive constant $C_N(D)$ such that
\begin{equation}\label{eq:primed-modulus}
 M^{-h}|Z'_M(D)|\longrightarrow C_N(D).
\end{equation}
Under \eqref{eq:J}, the primed Pfaffian has phase limit
$\exp(-\ii\pi\eta_D(0)/4)$ and
$M^{-h/2}|P'_M(D)|\to C_N(D)^{1/2}$.
\end{proposition}

\begin{proof}
The phase follows from \eqref{eq:theta-heat} because the heat trace
$T(t)$ is unchanged by removal of the kernel. For the modulus, let
\[
 S(t)=\sum_{k=0}^r a_{2k}t^{k-n/2},\qquad
 R(t)=K(t)-h-S(t)+h\mathbf1_{(0,1)}(t).
\]
Then $R(t)/t$ is absolutely integrable: at zero $R(t)=O(t^{1/2})$,
and at infinity $K(t)-h$ decays exponentially while $S(t)$ decays
as negative powers. The same finite-mass calculation and
Lemma~\ref{lem:heat-moments} give
\begin{equation}\label{eq:primed-log}
 \log|Z'_M(D)|
 =-\frac12\int_0^\infty B_M(t)R(t)\frac{\dd t}{t}
   +\frac h2\int_0^1 B_M(t)\frac{\dd t}{t}.
\end{equation}
The first integral has a finite limit by dominated convergence.
Write $b(u)=\sum_{j=0}^Nc_je^{-j^2u}$, so that $B_M(t)=b(M^2t)$.
Since $b(u)=O(u^{N/2})$ at zero and $b(u)-1$ decays exponentially
at infinity,
\begin{equation}\label{eq:b-constant}
 \int_0^1 B_M(t)\frac{\dd t}{t}-2\log M
 \longrightarrow
 \beta_N:=\int_0^1\frac{b(u)}u\dd u
          +\int_1^\infty\frac{b(u)-1}{u}\dd u.
\end{equation}
Thus \eqref{eq:primed-modulus} holds with
\begin{equation}\label{eq:C-primed}
 C_N(D)=\exp\!\left[-\frac12\int_0^\infty R(t)\frac{\dd t}{t}
                              +\frac h2\beta_N\right].
\end{equation}
The Pfaffian assertion follows by the eigenspace construction of
Theorem~\ref{thm:pf}. A fixed system of spectral and mass units is
understood in these normalization formulas.
\end{proof}

In particular, $\log|Z'_M|=h\log M+O(1)$ and
$\log|P'_M|=(h/2)\log M+O(1)$. Removing zero modes leaves the eta
phase intact, but changes the normalization needed for the modulus.

\section{Extensions and scope}\label{sec:extensions}

\subsection{An abstract heat-trace criterion}

The phase calculation can be separated from the odd-dimensional
geometry used to establish its analytic input.

\begin{proposition}\label{prop:abstract}
Let $D$ be a self-adjoint first-order elliptic operator on a closed
$n$-dimensional manifold. Suppose that it satisfies the Weyl
estimate \eqref{eq:weyl} and
\begin{equation}\label{eq:abstract-integrability}
 \int_0^\infty t^{-1/2}
            |\Tr(De^{-tD^2})|\dd t<\infty.
\end{equation}
Set
\begin{equation}\label{eq:eta-heat-definition}
 \eta_{\mathrm{heat}}(D)=\frac1{\sqrt\pi}
           \int_0^\infty t^{-1/2}\Tr(De^{-tD^2})\dd t.
\end{equation}
For any even $N>n$, the primed massless PV product satisfies
\begin{equation}\label{eq:abstract-phase}
 \lim_{M\to\infty}\ph Z'_M(D)
       =e^{-\ii\pi\eta_{\mathrm{heat}}(D)/2}.
\end{equation}
If the spectral eta function is regular at zero, then
$\eta_{\mathrm{heat}}(D)=\eta_D(0)$.
The negative-mass phase and the Pfaffian phases also hold with
$\eta_{\mathrm{heat}}(D)$ in place of $\eta_D(0)$, with
\eqref{eq:J} imposed for Pfaffians.
\end{proposition}

\begin{proof}
Finite-mass absolute convergence uses only the Weyl estimate and the
moment identities. The Fubini estimate \eqref{eq:fubini} is likewise
unchanged. Equation~\eqref{eq:abstract-integrability} supplies exactly
the majorant required in \eqref{eq:H-majorant}, proving the limit.
For $\Re s>0$, the Mellin integral \eqref{eq:mellin} is holomorphic
and agrees with the spectral eta function where its defining series
converges. Dominated convergence along real $s\downarrow0$ identifies
its limit with \eqref{eq:eta-heat-definition}. If eta is regular at
zero, this limit is its value there. The signed-mass kernel
\eqref{eq:H-negative} and the Pfaffian construction prove the
remaining assertions.
\end{proof}

The stronger estimate $T(t)=O(t^{1/2})$ directly implies regularity
as in Lemma~\ref{lem:eta}. The proposition does not extend the
odd-dimensional modulus formula without additional heat-kernel
information.

\subsection{Oriented even-dimensional manifolds}

For a nonchiral geometric Dirac operator on the full spinor bundle
of a closed oriented even-dimensional manifold, the chirality
operator $\Gamma$ satisfies $\Gamma D=-D\Gamma$. Its nonzero
spectrum is symmetric, so $T(t)=0$ and $\eta_D(0)=0$.
Moreover, \eqref{eq:ell} obeys
$\ell_M(-\lambda)=\overline{\ell_M(\lambda)}$.
Thus the primed massless PV determinant is positive at every $M$,
and its phase is already one before the large-mass limit.

The heat trace now has integer rather than half-integer powers.
The vanishing odd-moment argument for the modulus therefore does
not apply, and real local subtractions are generally needed.
For negative physical mass, Proposition~\ref{prop:abstract} gives
the additional zero-mode sign in \eqref{eq:negative-phase}.

\subsection{A covering-space criterion and Pin examples}

On a nonorientable manifold, global chirality need not exist.
Even-dimensional eta phases can consequently be nontrivial.
The operator and its Pin-module structure must first be specified.
The following sufficient criterion is useful without making a
universal assertion about every Pin Dirac convention.

Suppose $\widetilde Y\to Y$ is a double cover with a fixed-point-free
deck involution $\tau$. Let $\widetilde D$ be a lifted geometric
Dirac-type operator and let $U_\tau$ be a unitary involution lifting
$\tau$, commuting with $\widetilde D$. Assume that $D$ acts on the
invariant sections and that a unitary grading
$\widetilde\Gamma$ anticommutes with $\widetilde D$. Then
\begin{equation}\label{eq:cover-trace}
 T_Y(t)=\frac12\Tr(\widetilde D e^{-t\widetilde D^2})
       +\frac12\Tr(U_\tau\widetilde D e^{-t\widetilde D^2}).
\end{equation}
The first trace vanishes by the grading. Compactness and the absence
of fixed points give a positive lower bound on the distance from
$x$ to $\tau x$. The usual off-diagonal Gaussian estimate, also for
the differentiated heat kernel, bounds the second trace by
$O(t^{-L}e^{-c/t})$ for some $L$ and $c>0$.
This proves the short-time integrability required in
Proposition~\ref{prop:abstract}.

The four-dimensional $\mathrm{Pin}^+$ Majorana system discussed in
\cite[Appendices A and C]{Witten} provides the standard example of
a nontrivial even-dimensional eta phase. With the operator and
multiplicity conventions used there, the two $\mathrm{Pin}^+$
structures on the round $\R P^4$ give $h=0$ and
$\eta_D(0)=\pm1/4$. The negative-mass Pfaffian phase is accordingly
$e^{\mp\ii\pi/8}$. This example uses the specified Pin structure
and operator; parity of the dimension alone does not determine the
answer.

\subsection{Chiral fields, boundaries, and finite local terms}

A chiral operator $D^+:S^+\to S^-$ is not a self-adjoint operator
on a single space of sections. Its determinant, and the corresponding
Pfaffian in a real theory, generally takes values in a line bundle
over the background fields. Higher-dimensional eta invariants and
anomaly inflow enter its global interpretation
\cite{DF,FM,WY}. The scalar products constructed here cannot simply
be applied to $D^+$ by replacing $D$ in their formulas.

On a manifold with boundary, a self-adjoint domain and suitable heat
estimates must likewise be supplied. Our main theorem is for closed
manifolds. Finally, a finite imaginary local term in the effective
action can change a phase. All equalities in this paper refer to
the mass sign \eqref{eq:mass-sign} and the explicitly specified PV
normalization. Independence of $N$ in Theorem~\ref{thm:main} is an
independence within that family of schemes, not a statement about
arbitrary finite counterterms.

\section{Discussion}

The construction separates the ultraviolet cancellation from the
spectral asymmetry. Finite differences make the combined rational
spectral factor trace class at each regulator scale. The eta
invariant then enters through an absolutely convergent heat integral,
and the Bismut--Freed estimate justifies passage to infinite mass.
In odd dimensions the same regulator moments annihilate the
divergent heat terms in the modulus, yielding the conventional
positive zeta-determinant amplitude.

For Pfaffians, the antiunitary symmetry specifies finite-dimensional
antisymmetric forms on the eigenspaces. Their PV ratios fix the
square root, including its spectral-flow sign. Negative physical
mass changes the limiting heat kernel by a factor of two, while
zero modes contribute the explicit signs in
Theorem~\ref{thm:negative}. In a massless primed determinant they
instead change the power of the regulator scale required to
normalize the modulus.

The finite-product convergence proofs give an explicit link between
the physical eta phases of \cite{Witten,WY} and the heat-kernel
results of \cite{BF,Ponge}. Their hypotheses also make precise the
distinction from the determinant-line constructions of \cite{DF,FM}.

\clearpage
\begin{thebibliography}{9}

\bibitem{Witten}
E.~Witten,
\emph{Fermion Path Integrals And Topological Phases},
Rev. Mod. Phys. \textbf{88} (2016), 035001,
\href{https://arxiv.org/abs/1508.04715}{arXiv:1508.04715}.

\bibitem{WY}
E.~Witten and K.~Yonekura,
\emph{Anomaly Inflow and the $\eta$-Invariant},
\href{https://arxiv.org/abs/1909.08775}{arXiv:1909.08775}.

\bibitem{BF}
J.-M.~Bismut and D.~S.~Freed,
\emph{The Analysis of Elliptic Families II. Dirac Operators,
Eta Invariants, and the Holonomy Theorem},
Commun. Math. Phys. \textbf{107} (1986), 103--163.
\href{https://webhomes.maths.ed.ac.uk/~v1ranick/papers/bismfree2.pdf}{Available PDF}.

\bibitem{Ponge}
R.~Ponge,
\emph{A new proof of the local regularity of the eta invariant of
a Dirac operator},
J. Geom. Phys. \textbf{56} (2006), 1654--1665,
\href{https://arxiv.org/abs/math/0411123}{arXiv:math/0411123}.

\bibitem{DF}
X.~Dai and D.~S.~Freed,
\emph{Eta-Invariants and Determinant Lines},
J. Math. Phys. \textbf{35} (1994), 5155--5194,
\href{https://arxiv.org/abs/hep-th/9405012}{arXiv:hep-th/9405012}.

\bibitem{FM}
D.~S.~Freed and G.~W.~Moore,
\emph{Setting the Quantum Integrand of M-Theory},
Commun. Math. Phys. \textbf{263} (2006), 89--132,
\href{https://arxiv.org/abs/hep-th/0409135}{arXiv:hep-th/0409135}.

\end{thebibliography}
\end{document}
